\par
(c) Sea $(S_0,\xi_0)$ un objeto formal mínimamente versal de
$\underline h_{(R_0,R_1,\ldots)}$ (Teorema 1.11). Para todo objeto
\[
 \xi:R_0\longrightarrow T
 \qquad\text{de }\underline h_{(R_0,R_1,\ldots)}(T),
\]
considérese el funtor
\[
 \operatorname{Isom}(\xi,\xi):
 \underline C_\Lambda/(T\widehat\otimes_\Lambda T)
 \longrightarrow(\mathrm{Sets}).
\]
Este asigna a $R$ el conjunto de isomorfismos entre las dos imágenes de
$\xi$ sobre $R$ obtenidas a partir de
\[
 R_0\xrightarrow{\ \xi\ }T
 \mathrel{\substack{\xrightarrow{}\\[-.6ex]\xrightarrow[]{}}}
 T\widehat\otimes_\Lambda T\longrightarrow R.
\]
Este funtor está representado por
\[
 R_1\otimes_{R_0\widehat\otimes_\Lambda R_0}
       (T\otimes_\Lambda T).
\]
El paso al límite inverso muestra que
$\operatorname{Isom}(\xi_0,\xi_0)$ está prorepresentado por un objeto
$S_1$, y que
\[
 \underline h_{(R_0,R_1,\ldots)}
 \simeq\underline h_{(S_0,S_1,\ldots)}.
\]
Por (b), el cogrupoide $S_*$ es mínimamente liso. \hfill$\square$

\underline{Definición 2.13.} Un grupoide cofibrado $\underline A$
sobre $\underline C_\Lambda$ se llama \underline{(mínimamente)
prorepresentable} si es $\underline C_\Lambda$-equivalente a
$\underline h_{(R_0,R_1,\ldots)}$ para algún cogrupoide (normalizado)
$(R_0,R_1,\ldots)$ en $\widehat{\underline C}_\Lambda$.

\underline{Proposición 2.14.} Si $(R_0,R_1,\ldots)$ es un cogrupoide
normalizado en $\widehat{\underline C}_\Lambda$, entonces todo
homomorfismo
\[
 (\alpha_0,\alpha_1):(R_0,R_1,\ldots)
 \longrightarrow(R_0,R_1,\ldots)
\]
que induce una $\underline C_\Lambda$-equivalencia
\[
 \underline h_{(\alpha_0,\alpha_1)}:
 \underline h_{(R_0,R_1,\ldots)}
 \longrightarrow\underline h_{(R_0,R_1,\ldots)}
\]
es un isomorfismo. En

%%%%%%%%%%  scan idx 77  |  folio 66  %%%%%%%%%%

particular, una prorepresentación normalizada de un grupoide cofibrado
sobre $\underline C_\Lambda$, si existe, es única salvo isomorfismo.

\underline{Demostración.} Considérese el diagrama conmutativo
\[
\begin{tikzcd}[column sep=huge,row sep=huge]
 \underline h_{(R_0,R_1,\ldots)}
   \arrow[r,"{\underline h_{(\alpha_0,\alpha_1)}}"]
 & \underline h_{(R_0,R_1,\ldots)}\\
 \underline h_{R_0} \arrow[u,"\epsilon"]
   \arrow[r,"{\underline h_{\alpha_0}}"']
 & \underline h_{R_0} \arrow[u,"\epsilon"']
\end{tikzcd}
\]
en el que $\underline h_{(\alpha_0,\alpha_1)}$ es una
$\underline C_\Lambda$-equivalencia. Puesto que
$\underline h_{(R_0,R_1,\ldots)}(k[\varepsilon])$ es totalmente
disconexo, tanto
\[
 [\underline h_{(\alpha_0,\alpha_1)}(k[\varepsilon])]
 \quad\text{como}\quad
 [\epsilon(k[\varepsilon])]
\]
son biyecciones, y por tanto también lo es
$\underline h_{\alpha_0}(k[\varepsilon])$. Así,
$\alpha_0:R_0\longrightarrow R_0$ es un endomorfismo suprayectivo y,
por consiguiente, un automorfismo. Se sigue que $\alpha_1$ es también un
automorfismo de $R_1$. \hfill$\square$
